% Section 5 — Learning from verified outcomes
\section{Learning from verified outcomes}
\label{sec:training}

EurekaZero learns the way AlphaGo Zero did, with one change of venue (\cref{fig:loop}): during
training, search plays the innovation games and an external verifier scores every move; the network
learns to make the first move; at test time the network moves alone. The improvement operator is
search plus verification; the network's job is to carry what search found into a single proposal on
problems search has not yet seen \citep{silver2017zero,anthony2017exit}.

\begin{figure}[t]
\centering
% Figure — the discovery loop: search plays the games during training; the network learns to make
% the first move. Conceptual; no data.
\begin{tikzpicture}[x=1cm, y=1cm,
  b/.style={ezbox, font=\sffamily\scriptsize, text width=2.15cm, minimum height=1.05cm, inner sep=3pt},
  s/.style={ezstage, font=\sffamily\scriptsize, text width=2.15cm, minimum height=1.05cm, inner sep=3pt},
  lab/.style={font=\sffamily\fontsize{5.8pt}{6.6pt}\selectfont, text=ezgrey, align=center, inner sep=1.5pt}]
  % training half
  \fill[ezpaper, rounded corners=4pt] (-0.2,-2.8) rectangle (9.65,2.35);
  \node[font=\sffamily\sbseries\scriptsize, text=ezgrey, anchor=south west] at (-0.1,-2.75) {TRAINING: search plays the games};
  \node[b] (fam) at (1.3,0.75) {\textbf{Training families}\\states $s$ with\\incumbents $h(s)$};
  \node[ezhidden, font=\sffamily\scriptsize, text width=2.15cm, minimum height=1.05cm, inner sep=3pt] (srch) at (4.75,0.75)
    {\textbf{Verified search}\\teacher at budget $k$\\{\color{white!70}protected}};
  \node[b] (store) at (8.2,0.75) {\textbf{Experience store}\\every attempt,\\failures kept};
  \node[s] (stg) at (8.2,-1.25) {\textbf{Stages A--D}\\validity \textperiodcentered{} imitation\\$J_1$ \textperiodcentered{} rounds};
  \node[s] (net) at (4.75,-1.25) {\textbf{EurekaZero}\\$P_\theta$, $D_{\psi,d}$};
  \draw[ezarrow] (fam) -- (srch);
  \draw[ezarrow] (srch) -- node[lab, above] {verified\\moves} (store);
  \draw[ezarrow] (store) -- (stg);
  \draw[ezarrow] (stg) -- node[lab, above] {weights} (net);
  \draw[ezarrow, dashed] (net.west) -| node[lab, pos=0.3, below] {seeds the next round} (fam.south);
  \draw[ezarrow, dashed] (srch.north) to[out=110, in=70, looseness=0.75]
    node[lab, above] {wins raise the incumbent $h(s)$} (fam.north);

  % deployment half
  \fill[ezaccentsoft, rounded corners=4pt] (9.9,-2.8) rectangle (15.1,2.35);
  \node[font=\sffamily\sbseries\scriptsize, text=ezaccent, anchor=north west] at (10.0,2.3) {TEST TIME: Zero};
  \node[b, text width=2.3cm] (new) at (12.6,1.05) {\textbf{Unseen state} $s$\\new dataset, history\\or specification};
  \node[s, text width=2.3cm] (move) at (12.6,-0.35) {\textbf{One move} $a$\\committed, hashed};
  \node[b, text width=2.3cm] (ver) at (12.6,-1.75) {\textbf{Verifier}\\outcome $u$\\after commitment};
  \draw[ezarrow] (new) -- (move);
  \draw[ezarrow] (move) -- (ver);
  \draw[ezarrow, ezaccent] (net.south) -- (4.75,-2.35) -- node[lab, above, pos=0.4, text=ezaccent] {frozen weights}
    (10.35,-2.35) |- (move.west);
  \node[lab, text=ezwarm, anchor=south] at (12.6,-2.75) {no candidate tested before commitment};
\end{tikzpicture}

\caption{The discovery loop. In training, a verified search plays each training game at budget $k$;
every attempt, including failures, enters the experience store; staged training produces EurekaZero,
whose proposals seed the next round, and every verified win raises that state's incumbent. At test
time the frozen network makes one move on an unseen state, and the verifier scores it only after
commitment. The search's internals are protected; its budget, outputs and receipts are not.}
\label{fig:loop}
\end{figure}

\subsection{The unit of experience}

A training example binds the state, the candidate, the verifier's result, its cost and its
provenance, and the store keeps the full attempted set---inferior candidates and failures
included---rather than only winners (\cref{app:corpus}). Each record names the policy that produced
the candidate; a selection probability is recorded only when it was actually known, never inferred
from an archive rank, because unknown propensities limit every off-policy claim.

\subsection{Teachers}

The principal teacher is a verified search over the solution language, run at a declared budget on
training families. Conventional solvers, qualified human artifacts and verified candidates from earlier
EurekaZero rounds can join the corpus, each with its source recorded, and a \emph{search-only} arm
trained without them measures what they add (\cref{sec:mechanisms}). The corpus and its family map
are frozen before any confirmation family is exposed, records from a teacher that saw a family's test
outcomes are excluded from claims about that family, and examples are balanced by family so that a few
cheap tasks with millions of candidates cannot dominate.

\subsection{Four training stages}

\paragraph{Stage A: representation and validity} The decoders learn to reconstruct decomposed
training artifacts and to predict feasibility. We measure reconstruction and validity on held-out
artifact families, and claim nothing about innovation from this stage.

\paragraph{Stage B: conditional imitation} The proposer learns $P_\theta(h\mid z)$ from states and
verified teacher outcomes, weighting candidates by verified utility and keeping several strong,
structurally different solutions per state. Uniform imitation and strongest-solution imitation are
compared arms.

\paragraph{Stage C: the first-proposal objective} The objective that matches the Zero contract is the
expected utility of the one move the deployed policy will make,
\begin{equation}
J_1(\theta)=\mathbb E_{\tau\sim p_{\mathrm{train}}}\ \mathbb E_{a\sim\pi_\theta(\cdot\mid s_\tau)}\bigl[U_\tau(a)\bigr],
\label{eq:j1}
\end{equation}
estimated on policy with a per-state baseline $v(s)$ and a regularizer toward a reference policy,
\begin{equation}
\mathcal L_{\mathrm{PG}}=-\frac1m\sum_{i=1}^{m}\bigl(u_i-\operatorname{sg}(v(s_i))\bigr)\log\pi_\theta(a_i\mid s_i)
+\beta\,\mathrm{KL}\bigl(\pi_\theta\,\Vert\,\pi_{\mathrm{ref}}\bigr).
\label{eq:pg}
\end{equation}
The reference must give positive probability to every legal move: a reference whose support is a
finite catalog makes the KL term infinite for every new solution and silently forbids invention. We
therefore use a declared mixture of the Stage-B policy and a grammar-uniform policy. Checkpoints are
selected on $J_1$ at $K=1$ on development families, never on best-of-$B$, which can improve while the
first proposal gets worse; and a falling training loss is not evidence of a better first proposal,
a lesson the program learned the hard way (\cref{app:history}).

\paragraph{Stage D: controlled improvement rounds} New training-family candidates are generated,
verified and added; the model is updated; and a development panel reserved for this stage measures
the first proposal. Each round is compared with adding the same amount of data from a static teacher
and with training the same model longer on the old data, and confirmation families never enter the
loop. This is the stage that turns imitation into self-play against the incumbent.

\subsection{Sequential games}

For the next-experiment game, a useful exploratory query may not improve the incumbent immediately.
We therefore train two objectives and report them separately: the one-step utility of the next
decision given a frozen history (track CH), and the final regret after $T$ decisions (track SEQ),
optimized as a horizon-dependent return. The teacher for both is EurekaZero's own look-ahead search on
training problems, where the objective is cheap to query; imitating the traces of human-designed
optimizers is a labelled baseline arm, not the training signal.

\Cref{tab:training} gives the training ledger. Its costs reappear, uncompressed, in the break-even
analysis of \cref{sec:compute}.

\begin{table}[t]
\caption{Training ledger by domain. All entries are synthetic targets (gaps EZ-G18, EZ-G20).}
\label{tab:training}
\centering\tablesf
\begin{tabularx}{\linewidth}{@{}lLLL@{}}
\toprule
 & Learning procedures & Next experiment & CAD parts \\
\midrule
Teacher search (CPU-hours) & \syn[EZ-G18]{\EZteachA} & \syn[EZ-G18]{\EZteachB} & \syn[EZ-G18]{\EZteachC} \\
Confirmation families & \syn[EZ-G20]{\EZnA} & \syn[EZ-G20]{\EZnB} & \syn[EZ-G20]{\EZnC} \\
Proposal latency, median (s) & \syn[EZ-G18]{\EZlatA} & \syn[EZ-G18]{\EZlatB} & \syn[EZ-G18]{\EZlatC} \\
Training compute, all stages & \multicolumn{3}{l}{\syn[EZ-G18]{\EZgpuH} GPU-hours across the three domains} \\
\bottomrule
\end{tabularx}
\end{table}
